\section{Background and related work}
\label{sec:background}

\subsection{Leakage in bearing-diagnosis evaluation}

Hendriks, Dumond and Knox showed that the conventional condition-wise split of the CWRU data places the same physical
bearings in training and test sets, and proposed benchmarking with independent sets of bearings \cite{hendriks2022}. Wheat et
al.\ compared run-to-run, day-to-day and part-to-part splits for classical pipelines and found the gap between them large
enough to change method rankings \cite{wheat2024leakage}. Vieira et al.\ generalised bearing-wise partitioning across four
datasets, reformulated diagnosis as multi-label detection and measured the inflation caused by bearing- and
segmentation-level leakage over 100 evaluation splits \cite{vieira2026}; their study is supervised, and cross-domain and
source-free settings are stated as out of scope. Matania et al.\ \cite{matania2024leakage} make the same case for
condition-based maintenance generally, showing how test--training separation should be constructed. Knap et al.\ published a
leakage-safe cross-domain benchmark on CWRU and Paderborn with six fixed source--target scenarios
\cite{knap2026leakagesafe}; separation there is enforced at
\emph{recording} level, which removes window leakage but still permits the same physical bearing on both sides of a transfer
task, and no source-free method is evaluated.

Beyond this field, Kapoor and Narayanan \cite{kapoor2023leakage} document leakage as a systemic failure across 17
scientific disciplines and give the taxonomy we adopt, and Geirhos et al.\ \cite{geirhos2020shortcut} name the underlying
learning behaviour: a shortcut that succeeds on the benchmark and fails under the intended shift.

This paper takes the bearing-wise principle from that line of work and applies it where it has not been applied: to
source-free adaptation, with the source model shared between the leaky and the clean target so that the two differ in
bearing overlap and nothing else, and with the split treated as the statistical unit.

\subsection{Source-free domain adaptation for bearings}

SHOT established the template on which most later methods build: freeze the source classifier, maximise information at the
target and refine centroid-based pseudo-labels \cite{liang2020shot}. SDALR extends it by separating reliable from unreliable
target predictions, using the reliable ones as supervision and the unreliable ones through an entropy term, and reports
\SI{96.8}{\percent} mean accuracy across Paderborn operating conditions \cite{sdalr2025}. Jeong et al.\ combine
speed-normalised order spectra with a variational autoencoder and test-time training, and find that order normalisation alone
accounts for most of their target gain \cite{jeong2025}. Published audits of SFDA stability in this field address other
failure modes, such as collapse under industrial noise, rather than bearing overlap. In all of this work the target domain is
a different operating condition of the same physical bearings.

\subsection{Physics priors inside adaptation}

Characteristic frequencies computed from geometry and speed are the natural prior for bearing diagnosis, and recent
physics-guided methods use them to steer adaptation. PCTL scores predicted classes against envelope-spectrum energy at
characteristic frequencies and trains a confidence predictor, using a small set of labelled target data \cite{pctl2026}.
Bio-SFDA gates and weights pseudo-labels with band-energy and signal-to-noise predicates at BPFO, BPFI and BSF computed on a
short-time Fourier transform, with thresholds adapted online from unlabelled target statistics \cite{yoon2026biosfda}. These
modules are evaluated through end accuracy; the rate at which they accept a fault label on a healthy bearing is not reported.
Because the Bio-SFDA task definition for Paderborn cannot be reconstructed from the article --- a ball class is listed, and
the Paderborn real-damage set contains no rolling-element damage, and the split is not described --- we evaluate the gating
\emph{mechanism as described, as we implemented it}, and make no comparison against the accuracy reported there.

\subsection{Envelope analysis and the statistics of detection}

The squared envelope spectrum is the standard carrier of second-order cyclostationary bearing signatures \cite{randall2011},
usually after a band selection such as the fast kurtogram \cite{antoni2007kurtogram} and, where discrete components
dominate, after discrete/random separation. Borghesani et al.\ derived the statistics of the squared envelope spectrum under
coloured noise and the corresponding thresholds for testing cyclostationarity \cite{borghesani2013}. The gates compared here
are built from those ingredients; none of them is claimed as new. Matania et al.\ project order-domain envelope spectra onto
kinematic harmonic positions to obtain a machine-invariant representation \cite{matania2025}, which is the same idea as the
harmonic coordinates used internally by our comb tests, and is cited as such rather than claimed.

\subsection{Invariance to a nuisance that carries the label}

Removing a nuisance factor from a representation by adversarial training is standard, but when the nuisance carries label
information, unconditional invariance conflicts with classification: aligning the marginal feature distribution under label
shift induces concept shift and, at the invariant optimum, mixes the classes \cite{liu2021labelshift}. Class-conditional
adversaries restore compatibility by removing the nuisance \emph{within} each class \cite{li2018ciddg}, and
subject-adversarial training plays the same role in EEG decoding \cite{ozdenizci2020eeg}. A bearing source set is the extreme
case of the conflict, because each physical bearing carries exactly one fault class, so class is a deterministic function of
bearing identity. Section~\ref{sec:res-probe} measures how strongly that identity is present; the corresponding remedy is
left as future work and is not claimed here.
